Large Language Models (LLMs) show significant potential in economic and strategic interactions, where communication via natural language is often prevalent. This raises key questions: Do LLMs behave rationally?  How do they perform compared to humans? Do they tend to reach an efficient and fair outcome? What is the role of natural language in strategic interaction? How do characteristics of the economic environment influence these dynamics? These questions become crucial concerning the economic and societal implications of integrating LLM-based agents into real-world data-driven systems, such as online retail platforms and recommender systems.
While the ML community has been exploring the potential of LLMs in such multi-agent setups, varying assumptions, design choices and evaluation criteria across studies make it difficult to draw robust and meaningful conclusions. To address this, we introduce a benchmark for standardizing research on two-player, sequential, language-based games. Inspired by the economic literature, we define three base families of games with consistent parameterization, degrees of freedom and economic measures to evaluate agents' performance (self-gain), as well as the game outcome (efficiency and fairness). We develop an open-source framework for interaction simulation and analysis, and utilize it to collect a dataset of LLM vs. LLM interactions across numerous game configurations and an additional dataset of human vs. LLM interactions.
Through extensive experimentation, we demonstrate how our framework and dataset can be used to: (i) compare the behavior of LLM-based agents in various economic contexts; (ii) evaluate agents in both individual and collective performance measures; and (iii) quantify the effect of the economic characteristics of the environments on the behavior of agents. Our results suggest that the market parameters, as well as the choice of the LLMs, tend to have complex and interdependent effects on the economic outcome, which calls for careful design and analysis of the language-based economic ecosystem.